\ssr{ВВЕДЕНИЕ}

\textbf{Расстояние Левенштейна} -- минимальное количество операций вставки одного символа, удаления одного символа и замены одного символа на другой, необходимых для превращения одной строки в другую.

Оно используется для нечеткого поиска, коррекции опечаток, обработки текста и сравнения последовательностей, например, в биоинформатике и текстовых приложениях.

Целью данной лабораторной работы является проведение исследования метода динамического программирования на примере алгоритмов Левенштейна и Дамерау-Левенштейна. 

Задачами данной лабораторной являются:
\begin{enumerate}
	\item Рассмотрение алгоритмов Левенштейна и Дамерау-Левенштейна нахождения расстояния между строками;
	\item Разработка рекурсивного алгоритма расстояния Левенштейна, а также применение методов динамического программирования для разработки алгоритмов расстояний Левенштейна и Дамерау-Левенштейна с использованием кеширования; 
	\item Получение практических навыков реализации указанных алгоритмов: двух алгоритмов с кешированием и одного из алгоритмов в рекурсивной версии; 
	\item Экспериментальное подтверждение различий во временной эффективности рекурсивной и итерациной реализаций алгоритма определения расстояния Левенштейна с помощью замеров процессорного времени и пиковой памяти при выполнения алгоритмов на варьирующихся длинах строк.
\end{enumerate}
\newpage

\clearpage
